\section{The complete algorithm and certified outputs}\label{sec:completion}

The preceding results now fit together without a further probabilistic
argument. One root call recovers the potential, and one threshold test
per vertex converts it to the required set.

\begin{proof}[Proof of Theorem~\ref{thm:main}]
From the explicit input compute $t,W,B_*$ according to
Equation~\eqref{eq:transformed}, and initialize the global parameters
of Equation~\eqref{eq:global-parameters}. Run
\[
 p=\Solve(0,B_*;B_*,\boldsymbol1,\boldsymbol1,n,n)
\]
and return $U_p=\{i\in V:p_i>B_*/2\}$. Scalars in the box positions
denote constant vectors.

Lemma~\ref{lem:box-solution} gives the unique deterministic solution
$x$ of $[0,B_*]$. Its two supports have size at most $n$ and so satisfy
the root budgets. By Corollary~\ref{cor:top-success}, the single event
$p=x$ has probability at least $7/8$. On that event,
Proposition~\ref{prop:game-reduction} identifies $U_p$ with the entire
zero-threshold winning set for the original weights and liminf payoff,
including against every opposing history-dependent strategy.

The computation terminates on every random tape by
Proposition~\ref{prop:termination}. Its full bit cost is bounded by
Theorem~\ref{thm:bit-complexity}, which already includes the
preprocessing and the $n$ final threshold comparisons. The procedures
use uniform finite instructions on binary integers, exact rationals,
and fair bits. Thus they supply both the stated output probability and
the asserted bound on every tape.
\end{proof}


The potential also provides a certificate that can be checked without
repeating the randomized proof. This gives a version that reports failure
instead of ever returning an incorrect answer.

\begin{corollary}[Certified winning regions and strategies]
\label{cor:certified-output}
There is a randomized algorithm which, on every tape, uses
$2^{O((\log_2(L+2))^2)}$ bit operations and returns either a failure
symbol or the exact winning set together with an integer potential and
positional strategies for both players. Whenever it returns these
objects, the Max strategy wins from every vertex of the winning set and
the Min strategy ensures upper limiting average at most $-1/n$ from
every other vertex, against arbitrary opposing strategies. It returns
these objects with probability at least $7/8$.

Repeating independent trials until a certificate is accepted terminates
almost surely, never returns an incorrect answer, and has expected bit
cost $2^{O((\log_2(L+2))^2)}$.
\end{corollary}

\begin{proof}
Run the root solver to obtain $p\in[0,B_*]$ and check
\[
                 p=\clip_{[0,B_*]}F(p).
\]
Accept precisely when this equality holds. By
Lemma~\ref{lem:box-solution}, it holds if and only if $p=x$, the
unique top-box solution. Thus acceptance certifies every coordinate,
and Corollary~\ref{cor:top-success} gives acceptance probability at
least $7/8$.

On acceptance, form $U=\{i:p_i>B_*/2\}$. At each Max vertex in
$U$, choose an outgoing edge attaining $F_i(p)$; at each Min vertex
outside $U$, choose an attaining edge as well. Complete the two
strategies arbitrarily at their other vertices. The two-band proof of
Proposition~\ref{prop:game-reduction} applies to the verified vector
$p=x$ and proves both strategy guarantees. The check and the edge
selections require a scan of the explicit edge list and arithmetic on
$O(L+1)$-bit integers, so they add only polynomial bit cost to the
bounded-time trial.

For independent repetitions, the probability that the first $r$ trials
all reject is at most $8^{-r}$. Hence acceptance occurs almost surely,
and the expected number of trials is at most
$\sum_{r\ge0}8^{-r}=8/7$. Multiplying by the worst-case bit cost of
one trial proves the expected-time bound. This repeated procedure has
no asserted time bound on every tape; each individual certifying trial
retains that bound.
\end{proof}
